\SourcePage{335}{340}
\section{Photon propagator}

Until now it was necessary (in~\S\S\,43, 74) to use the explicit form of the operators of the electromagnetic field $\hat{A}$ when finding matrix elements only for transitions involving a change of the number of real photons. For this purpose it was sufficient to write the potentials of the free field in the form of the expansion~(\S\,2) in transverse plane waves.

Such a representation, however, does not by itself provide a complete description of the arbitrary field. This is already clear from the fact that the diagrams of scattering~(73.13), (73.14) should describe also the Coulomb interaction of the electrons. The latter is described by the scalar potential $\Phi$, and cannot obviously be reduced to the exchange of only transverse virtual photons.

These difficulties, however, have only a formal, not physical, character, and they can be circumvented by using some general properties of the propagator, obvious from the requirements of relativistic and gauge invariance.

The most general form of a tensor of second rank, depending only on the 4-vector $\xi = x - x'$, is
\begin{equation}
D_{\mu\nu}(\xi) = g_{\mu\nu}D(\xi^2) - \partial_\mu\partial_\nu D^{(l)}(\xi^2),
\tag{76.2}
\end{equation}
where $D$, $D^{(l)}$ are scalar functions of the invariant $\xi^2$\,\footnote{These functions are different in three regions of values of the argument, not transforming into one another under Lorentz transformations: outside the light cone ($\xi^2 < 0$), in the upper ($\xi^2 > 0$, $\xi_0 > 0$) and in the lower ($\xi^2 > 0$, $\xi_0 < 0$) halves of the light cone.}. Let us note that the tensor automatically turns out to be symmetric. Correspondingly in the momentum representation we shall have
\begin{equation}
D_{\mu\nu}(k) = D(k^2)\,g_{\mu\nu} + k_\mu k_\nu D^{(l)}(k^2),
\tag{76.3}
\end{equation}
where $D(k^2)$, $D^{(l)}(k^2)$ are the Fourier components of $D(\xi^2)$, $D^{(l)}(\xi^2)$.

In the physical quantities---the amplitudes of scattering---the photon propagator enters multiplied by the transition currents of two electrons, i.e.\ in combinations of the form $j^\mu_{21}D_{\mu\nu}j^\nu_{43}$ (see, for example,~(73.13)). But by virtue of conservation of the current ($\partial_\mu j^\mu = 0$) its matrix elements $j_{21} = \bar{\psi}_2\gamma\psi_1$ satisfy the condition of 4-transversality
\begin{equation}
k_\mu(j^\mu)_{21} = 0,
\tag{76.4}
\end{equation}
where $k = p_2 - p_1$ (cf.~(43.13)). It is clear therefore that the physical results will not change under the replacement
\begin{equation}
D_{\mu\nu} \to D_{\mu\nu} + \chi_\mu k_\nu + \chi_\nu k_\mu,
\tag{76.5}
\end{equation}

\SourcePage{336}{341}
where $\chi_\mu$ are any functions of $\mathbf{k}$ and $k_0$. This arbitrariness in the choice of $D_{\mu\nu}$ corresponds to the arbitrariness of gauge of the potentials of the field. An arbitrary gauge transformation~(76.5) may violate the relativistically invariant form of $D_{\mu\nu}$ assumed in~(76.3) (if the quantities $\chi_\mu$ do not constitute a 4-vector). But also remaining within the framework of relativistically invariant forms of the propagator, we see that the choice of the function $D^{(l)}(k^2)$ in~(76.3) is entirely arbitrary; it does not affect the physical results and can be specified from considerations of convenience (\textit{L.\,D.\,Landau, A.\,A.\,Abrikosov, I.\,M.\,Khalatnikov}, 1954).

The determination of the propagator reduces, thus, to the determination of one gauge-invariant

\SourcePage{337}{342}
function $D(k^2)$. If one considers a given value of $k^2$ and chooses the $z$-axis along the direction of $\mathbf{k}$, then the transformations~(76.5) will not affect the components $D_{xx} = D_{yy} = -D(k^2)$. It is therefore sufficient to compute just one component $D_{xx}$.

Let us use the gauge in which $\mathrm{div}\,\hat{\mathbf{A}} = 0$ and the operator $\hat{A}$ is given by the expansion~(2.17), (2.18):
\begin{equation}
\hat{A} = \sum_{\mathbf{k}\alpha}\sqrt{\frac{2\pi}{\omega}}\,(\hat{c}_{\mathbf{k}\alpha}\,\mathbf{e}^{(\alpha)}e^{-ikx} + \hat{c}_{\mathbf{k}\alpha}^+\,\mathbf{e}^{(\alpha)*}e^{ikx}),\quad \omega = |\mathbf{k}|
\tag{76.6}
\end{equation}
(the index $\alpha = 1, 2$ labels the polarisations). From all vacuum averages of products of creation and annihilation operators only the following is different from zero:
\[
\langle 0|c_{\mathbf{k}\alpha}\,c_{\mathbf{k}\alpha}^+|0\rangle = 1.
\]

From definition~(76.1) we obtain therefore
\begin{equation}
D_{ik}(\xi) = \frac{1}{(2\pi)^3}\int\frac{2\pi i\,d^3k}{\omega}\left(\sum_\alpha e_i^{(\alpha)}e_k^{(\alpha)*}\right)e^{-i\omega|\tau|+i\mathbf{k}\boldsymbol{\xi}}
\tag{76.7}
\end{equation}
($i$, $k$ are three-dimensional vector indices; from summation over $\mathbf{k}$ we have passed to integration over $d^3k/(2\pi)^3$). The fact that in the exponent stands the absolute value of the difference $\tau = t - t'$ is a consequence of the chronologisation of the product of operators in~(76.1).

From~(76.7) it is seen that the sub-integral expression without the factor $e^{i\mathbf{k}\boldsymbol{\xi}}$ is the Fourier component of the three-dimensional function $D_{ik}(\mathbf{r},t)$. For $D_{xx} = -D$ it equals
\[
\frac{2\pi i}{\omega}\,e^{-i\omega|\tau|}\sum_\alpha|e_x^{(\alpha)}|^2 = \frac{2\pi i}{\omega}\,e^{-i\omega|\tau|}.
\]

It remains to decompose $D_{xx}(k^2)$ in the Fourier transform with respect to time:
\[
\frac{2\pi i}{\omega}\,e^{-i\omega|\tau|} = -\frac{1}{2\pi}\int\frac{4\pi}{k_0^2 - \mathbf{k}^2 + i0}\,e^{-ik_0\tau}\,dk_0.
\]

As was explained in the preceding section, such an integration implies going around the poles $k_0 = \pm|\mathbf{k}| = \pm\omega$ respectively from below and above; when $\tau > 0$ the integral is determined by the residue in the pole $k_0 = +\omega$, and when $\tau < 0$ by the residue in the pole $k_0 = -\omega$. In both cases the required result is obtained.

Thus we find finally
\begin{equation}
D(k^2) = \frac{4\pi}{k^2 + i0}.
\tag{76.8}
\end{equation}

\SourcePage{338}{343}
The appearance of $+i0$ in the denominator, to which in the above derivation we arrived automatically, coincides with the rule~(75.15): from (the zero) mass of the photon $i0$ is subtracted. From~(76.8) it is seen that the corresponding coordinate function $D(\xi^2)$ satisfies the equation
\begin{equation}
-\partial_i\partial^i D(x - x') = 4\pi\delta^{(4)}(x - x'),
\tag{76.9}
\end{equation}
i.e.\ it is the Green's function of the wave equation.

We shall ordinarily use the propagator in the form
\begin{equation}
D_{\mu\nu} = g_{\mu\nu}\,D(k^2) = \frac{4\pi}{k^2 + i0}\,g_{\mu\nu}
\tag{76.10}
\end{equation}
(\textit{Feynman gauge}).

Let us indicate also other gauges, which can present certain advantages in some applications.

Setting $D^{(l)} = -D/k^2$, we obtain the propagator in the form
\begin{equation}
D_{\mu\nu} = \frac{4\pi}{k^2}\left(g_{\mu\nu} - \frac{k_\mu k_\nu}{k^2}\right)
\tag{76.11}
\end{equation}
(\textit{Landau gauge}). In this case $D_{\mu\nu}k^\nu = 0$. Such a choice is analogous to the Lorentz gauge of the potentials ($A_\mu k^\mu = 0$).

To the gauge of the potentials by the three-dimensional condition $\mathrm{div}\,\mathbf{A} = 0$ the analogous gauge of the propagator is defined by the conditions
\[
D_{il}k^l = 0,\qquad D_{0l}k^l = 0.
\]
Together with the equality $D_{xx} = -D = -4\pi/k^2$ these conditions give
\begin{equation}
D_{il} = \frac{4\pi}{\omega^2 - \mathbf{k}^2}\left(\delta_{il} - \frac{k_ik_l}{\mathbf{k}^2}\right).
\tag{76.12}
\end{equation}
In order to obtain also $D_{il}$, one must perform the transformation~(76.5), setting
\[
\chi_0 = -\frac{4\pi\omega}{2(\omega^2 - \mathbf{k}^2)\mathbf{k}^2},\qquad \chi_i = \frac{4\pi k_i}{2(\omega^2 - \mathbf{k}^2)\mathbf{k}^2}.
\]
In this case for the remaining components of $D_{\mu\nu}$ we obtain
\begin{equation}
D_{00} = -4\pi/\mathbf{k}^2,\qquad D_{0i} = 0.
\tag{76.13}
\end{equation}
Such a gauge is called the \textit{Coulomb gauge} (\textit{E.\,Salpeter}, 1952); let us note that $D_{00}$ here is the Fourier component of the Coulomb potential.

Finally, the gauge of the potentials by the condition $\Phi = 0$ is analogous to the gauge of the propagator in which
\begin{equation}
D_{il} = -\frac{4\pi}{\omega^2 - \mathbf{k}^2}\left(\delta_{il} - \frac{k_ik_l}{\omega^2}\right),\qquad D_{0i} = D_{00} = 0.
\tag{76.14}
\end{equation}

\SourcePage{339}{344}
This form is found convenient for applications in non-relativistic problems (\textit{I.\,E.\,Dzyaloshinskii, L.\,P.\,Pitaevskii}, 1959).

All the expressions given pertain to the momentum representation of the propagator. In some problems it is convenient to use the mixed frequency-coordinate representation, i.e.\ the function
\begin{equation}
D_{\mu\nu}(\omega,\mathbf{r}) = \int D_{\mu\nu}(\omega,\mathbf{k})\,e^{i\mathbf{k}\mathbf{r}}\,\frac{d^3k}{(2\pi)^3}.
\tag{76.15}
\end{equation}

In the Feynman gauge~(76.10)
\[
D_{\mu\nu}(\omega,\mathbf{r}) = g_{\mu\nu}D(\omega,\mathbf{r}),
\]
where
\[
D(\omega,\mathbf{r}) = 4\pi\int\frac{e^{i\mathbf{k}\mathbf{r}}}{(2\pi)^3}\,\frac{d^3\mathbf{k}}{\omega^2 - \mathbf{k}^2 + i0} = -\frac{i}{\pi r}\int_0^\infty\frac{e^{ikr} - e^{-ikr}}{\omega^2 - k^2 + i0}\,k\,dk,
\]
or, after the replacement $k \to -k$ in the second term of the sub-integral expression:
\[
D(\omega,\mathbf{r}) = -\frac{i}{\pi r}\int_{-\infty}^{\infty}\frac{e^{ikr}\,k\,dk}{\omega^2 - k^2 + i0}.
\]

The last integration is performed by closing the contour of integration by an infinitely distant semicircle in the upper half-plane of the complex variable $k$ and reducing to the residue at the pole $k = |\omega| + i0$. We finally obtain
\begin{equation}
D(\omega,\mathbf{r}) = -e^{i|\omega|r}/r.
\tag{76.16}
\end{equation}

In connection with this expression let us make the following remark. The process described by diagrams~(73.13), (73.14) can be regarded as scattering of electron~2 in the field created by electron~1 (or vice versa). The function~(76.16) corresponds to the ordinary ``retarded'' potential $\alpha\,e^{i\omega r}$ (see~II, (64.1), (64.2)) only for $\omega > 0$. The sign $\omega$, however, depends on the conventional choice of direction of the arrow $k$ on the diagram. The noted property of the function $D(\omega,\mathbf{r})$ means that in quantum electrodynamics the source of the field should be considered that particle which emits energy, i.e.\ emits a virtual photon.

In conclusion let us dwell on the question of the propagator of a particle with spin~1 and with a mass different from zero. In this case the gauge arbitrariness is absent and the choice of the propagator is unambiguous.

Substituting the $\psi$-operators~(14.16) into the definition
\begin{equation}
G_{\mu\nu} = -i\langle 0|\mathrm{T}\psi_\mu(x)\psi^+_\nu(x')|0\rangle,
\tag{76.17}
\end{equation}

\SourcePage{340}{345}
we obtain an expression differing from~(76.7) only by the replacement of the sum over polarisations in the sub-integral expression by
\[
\sum_\alpha u^{(\alpha)}_\mu u^{(\alpha)*}_\nu.
\]

The summation over polarisations is equivalent, with a subsequent multiplication by~3, to averaging over the number of independent polarisations. The averaging gives the density matrix of unpolarised particles~(14.15). Thus we find the following expression for the propagator of massive vector particles:
\begin{equation}
G_{\mu\nu}(p) = -\frac{1}{p^2 - m^2 + i0}\left(g_{\mu\nu} - \frac{p_\mu p_\nu}{m^2}\right).
\tag{76.18}
\end{equation}

Let us call attention to the analogous structure of the propagators~(75.17) and~(76.18): in the denominator stands the difference $p^2 - m^2$, and the numerator is, to within a factor, the density matrix of unpolarised particles with a given spin.
